\documentclass[aspectratio=169,10pt]{beamer}

% Shared Beamer setup for Fall 2026 course slides.
\mode<presentation>{
  \usetheme{Madrid}
  \usecolortheme{default}
}

\definecolor{CalPolyGreen}{HTML}{154734}
\definecolor{CalPolyGold}{HTML}{BD8B13}
\definecolor{DeckInk}{HTML}{1F2933}
\definecolor{DeckMuted}{HTML}{5F6B7A}

\setbeamercolor{structure}{fg=CalPolyGreen}
\setbeamercolor{palette primary}{bg=CalPolyGreen,fg=white}
\setbeamercolor{palette secondary}{bg=DeckInk,fg=white}
\setbeamercolor{palette tertiary}{bg=CalPolyGold,fg=DeckInk}
\setbeamercolor{frametitle}{bg=CalPolyGreen,fg=white}
\setbeamercolor{title}{fg=CalPolyGreen}
\setbeamercolor{normal text}{fg=DeckInk,bg=white}
\setbeamercolor{block title}{bg=CalPolyGreen,fg=white}
\setbeamercolor{block body}{bg=black!3,fg=DeckInk}

\setbeamertemplate{navigation symbols}{}
\setbeamertemplate{footline}[frame number]
\setbeamertemplate{itemize item}{\color{CalPolyGold}$\blacktriangleright$}
\setbeamertemplate{itemize subitem}{\color{CalPolyGreen}$\bullet$}

\newcommand{\CourseNumber}{Course}
\newcommand{\CourseTitle}{Course Title}
\newcommand{\LectureNumber}{0}
\newcommand{\LectureTitle}{Lecture Title}
\newcommand{\LectureDate}{Date}
\newcommand{\CourseUnit}{Unit}

\newcommand{\coursenumber}[1]{\renewcommand{\CourseNumber}{#1}}
\newcommand{\coursetitle}[1]{\renewcommand{\CourseTitle}{#1}}
\newcommand{\lecturenumber}[1]{\renewcommand{\LectureNumber}{#1}}
\newcommand{\lecturetitle}[1]{\renewcommand{\LectureTitle}{#1}}
\newcommand{\lecturedate}[1]{\renewcommand{\LectureDate}{#1}}
\newcommand{\courseunit}[1]{\renewcommand{\CourseUnit}{#1}}

\author{Kirk Duran}
\institute{California Polytechnic State University, San Luis Obispo}
\date{\LectureDate}

\newcommand{\makedecktitle}{
  \begin{frame}[plain]
    \vspace{0.25in}
    {\usebeamerfont{title}\color{CalPolyGreen}\LectureTitle\par}
    \vspace{0.18in}
    {\large \CourseNumber{}: \CourseTitle\par}
    \vspace{0.08in}
    {\large Lecture \LectureNumber{} -- \LectureDate\par}
    \vfill
    {\small Kirk Duran \textbar{} Cal Poly, San Luis Obispo\par}
  \end{frame}
}

\newcommand{\placeholdernote}{
  \begin{block}{Placeholder}
    Replace these bullets with the final lecture material, examples, and in-class activities.
  \end{block}
}

\AtBeginSection[]{
  \begin{frame}{Roadmap}
    \tableofcontents[currentsection]
  \end{frame}
}


\usepackage{tikz}

\graphicspath{{../assets/lecture-09/}}

% If an image file is not yet available, skip the visual slot.
\newcommand{\lectureimage}[2][1.75in]{%
  \IfFileExists{../assets/lecture-09/#2}{%
    \includegraphics[width=\linewidth,height=#1,keepaspectratio]{#2}%
  }{%
  }%
}

% Temporary visible placeholder used while lecture assets are being assembled.
% Replace each placeholder with \lectureimage[...]{} once the corresponding asset exists.
\newcommand{\lectureplaceholder}[2][1.75in]{%
  \begin{center}
    \fbox{%
      \parbox[c][#1][c]{0.90\linewidth}{%
        \centering\small\textcolor{DeckMuted}{\textbf{Image placeholder}\\[0.08in]#2}%
      }%
    }
  \end{center}%
}

\setbeamersize{text margin left=0.42in,text margin right=0.42in}

\coursenumber{CS 4880}
\coursetitle{Artificial Intelligence}
\lecturenumber{9}
\lecturetitle{Constraint Satisfaction II}
\lecturedate{September 14, 2026}
\courseunit{Local and Constraint Search}

\begin{document}

\makedecktitle

\begin{frame}[shrink=20]{Today}
  \begin{itemize}
    \item Why correct backtracking can still explore an enormous search tree.
    \item How MRV, the degree heuristic, and LCV improve variable and value ordering.
    \item How forward checking detects immediate consequences of an assignment.
    \item How arc consistency and AC-3 propagate constraints across neighboring variables.
    \item How Maintaining Arc Consistency integrates propagation into recursive search.
    \item How backjumping and min-conflicts provide alternative ways to recover from or avoid bad decisions.
  \end{itemize}

  \begin{keyidea}
    Efficient CSP solving is not a different problem formulation; it is the systematic use of ordering and inference to avoid searching assignments that constraints have already ruled out.
  \end{keyidea}
\end{frame}

\sectionframe{Part 1: Why Naive Backtracking Is Not Enough}

\begin{frame}[shrink=20]{The Baseline Solver We Already Have}
  \begin{columns}[T,onlytextwidth]
    \begin{column}{0.55\textwidth}
      Lecture 8 gave us a complete recursive solver:
      \begin{enumerate}
        \item choose an unassigned variable;
        \item try a value from its domain;
        \item reject values inconsistent with the current assignment;
        \item recurse;
        \item undo the choice after failure.
      \end{enumerate}

      \begin{block}{Correctness is not the issue}
        With finite domains, basic backtracking is complete. The problem is how much of the search tree it may explore before finding a solution or proving failure.
      \end{block}
    \end{column}

    \begin{column}{0.39\textwidth}
      \lectureplaceholder[2.55in]{Backtracking search tree with a small successful branch hidden among many explored branches.}
    \end{column}
  \end{columns}
\end{frame}

\begin{frame}[shrink=20]{Where Does Backtracking Waste Time?}
  \begin{columns}[T,onlytextwidth]
    \begin{column}{0.30\textwidth}
      \begin{block}{Variable choice}
        We may assign an easy variable while a nearly impossible variable is waiting elsewhere.
      \end{block}
    \end{column}

    \begin{column}{0.30\textwidth}
      \begin{block}{Value choice}
        We may choose a legal value that destroys most of the remaining flexibility.
      \end{block}
    \end{column}

    \begin{column}{0.34\textwidth}
      \begin{block}{Late failure}
        We may continue searching long after the current assignments logically imply a contradiction.
      \end{block}
    \end{column}
  \end{columns}

  \vspace{0.12in}
  \begin{keyidea}
    The next improvements target three questions: where should we search, which value should we try, and what can we infer before searching further?
  \end{keyidea}
\end{frame}

\begin{frame}[shrink=20]{Three Levers for a Better CSP Solver}
  \begin{center}
  \renewcommand{\arraystretch}{1.35}
  \begin{tabular}{l|l|l}
    \textbf{Lever} & \textbf{Question} & \textbf{Techniques} \\\hline
    Variable ordering & Which variable next? & MRV, degree heuristic \\
    Value ordering & Which value first? & LCV \\
    Inference & What follows now? & Forward checking, AC-3, MAC
  \end{tabular}
  \end{center}

  \begin{keyidea}
    These techniques preserve the CSP itself; they change the order and amount of work needed to solve it.
  \end{keyidea}
\end{frame}

\sectionframe{Part 2: Variable and Value Ordering --- Search Where Failure Is Informative}

\begin{frame}[shrink=20]{Minimum Remaining Values (MRV)}
  \begin{cpdefinition}
    The \textbf{Minimum Remaining Values (MRV)} heuristic chooses the unassigned variable with the smallest number of currently legal values.
  \end{cpdefinition}

  \[
    X^*=\arg\min_{X\ \mathrm{unassigned}} |D_X|.
  \]

  \begin{itemize}
    \item A small domain means the variable is highly constrained by the current partial assignment.
    \item If a contradiction is inevitable, MRV tries to expose it early.
    \item MRV is therefore often called the \emph{fail-first} heuristic.
  \end{itemize}

  \begin{keyidea}
    Search the most fragile decision first rather than postponing it until after many irrelevant assignments.
  \end{keyidea}
\end{frame}

\begin{frame}[shrink=20]{Worked MRV Example}
  \begin{columns}[T,onlytextwidth]
    \begin{column}{0.50\textwidth}
      Suppose a partial map-coloring assignment leaves:
      \[
        D_{SA}=\{B\},
      \]
      \[
        D_Q=\{R,B\},
      \]
      \[
        D_{NSW}=\{R,G,B\}.
      \]

      MRV chooses $SA$ because
      \[
        |D_{SA}|=1.
      \]

      Assigning $SA$ now either succeeds immediately or exposes a contradiction immediately.
    \end{column}

    \begin{column}{0.44\textwidth}
      \lectureimage[2.60in]{slide9.png}
    \end{column}
  \end{columns}
\end{frame}

\begin{frame}[shrink=20]{Why MRV Can Shrink the Search Tree}
  \begin{columns}[T,onlytextwidth]
    \begin{column}{0.47\textwidth}
      \begin{block}{Arbitrary order}
        Choose a variable with four legal values, branch four ways, and discover much later that another variable had no valid continuation.
      \end{block}
    \end{column}

    \begin{column}{0.47\textwidth}
      \begin{block}{MRV order}
        Choose the nearly forced variable first. A dead end may be discovered before those four branches are ever created.
      \end{block}
    \end{column}
  \end{columns}

  \vspace{0.10in}
  \lectureimage[1.60in]{slide10.png}
\end{frame}

\begin{frame}[shrink=20]{Degree Heuristic}
  \begin{cpdefinition}
    The \textbf{degree heuristic} chooses the unassigned variable involved in the largest number of constraints on other unassigned variables.
  \end{cpdefinition}

  \begin{itemize}
    \item It is commonly used as a tie-breaker when several variables have the same MRV score.
    \item A high-degree variable can influence many remaining decisions.
    \item Assigning it early can create useful pruning elsewhere in the graph.
  \end{itemize}

  \begin{block}{Australia intuition}
    South Australia is highly connected, whereas Tasmania is isolated. If both have equally large domains, assigning South Australia is usually more informative.
  \end{block}
\end{frame}

\begin{frame}[shrink=20]{MRV and Degree Are Different Questions}
  \begin{center}
  \renewcommand{\arraystretch}{1.40}
  \begin{tabular}{l|l|l}
    \textbf{Heuristic} & \textbf{Looks at} & \textbf{Intuition} \\\hline
    MRV & current domain size & choose the most constrained variable \\
    Degree & remaining graph connectivity & choose the most constraining variable
  \end{tabular}
  \end{center}

  \begin{itemize}
    \item MRV is dynamic: domains may shrink as assignments and inference accumulate.
    \item Degree is also dynamic if we count only constraints involving still-unassigned variables.
    \item A common policy is MRV first, then degree as the tie-breaker.
  \end{itemize}

  \begin{keyidea}
    MRV asks where failure is most likely; degree asks where one decision can affect the largest portion of the remaining problem.
  \end{keyidea}
\end{frame}

\begin{frame}[shrink=20]{Least Constraining Value (LCV)}
  \begin{cpdefinition}
    The \textbf{Least Constraining Value (LCV)} heuristic orders the values of a chosen variable by how few values they eliminate from neighboring unassigned variables.
  \end{cpdefinition}

  \begin{itemize}
    \item Variable ordering tries to expose difficult variables early.
    \item Value ordering tries to preserve flexibility after that variable is chosen.
    \item LCV prefers values that leave neighbors with many alternatives.
  \end{itemize}

  \begin{keyidea}
    Choose the hardest variable first, then try the value that makes life easiest for the remaining variables.
  \end{keyidea}
\end{frame}

\begin{frame}[shrink=20]{Worked LCV Example}
  \begin{columns}[T,onlytextwidth]
    \begin{column}{0.53\textwidth}
      Suppose $Q$ has domain
      \[
        D_Q=\{R,G,B\}.
      \]

      Its unassigned neighbors have domains:
      \[
        D_{SA}=\{R,B\},\qquad D_{NSW}=\{R,G\}.
      \]

      Under inequality constraints:
      \begin{itemize}
        \item $Q=R$ removes one value from both neighbors;
        \item $Q=G$ removes no value from $SA$ and one from $NSW$;
        \item $Q=B$ removes one from $SA$ and none from $NSW$.
      \end{itemize}

      LCV prefers $G$ or $B$ before $R$.
    \end{column}

    \begin{column}{0.41\textwidth}
      \lectureimage[2.70in]{slide14.png}
    \end{column}
  \end{columns}
\end{frame}

\begin{frame}[shrink=20]{Ordering Heuristics Change the Search Tree}
  \begin{block}{Typical policy}
    \ttfamily\small
    X $\leftarrow$ unassigned variable chosen by MRV\\
    break MRV ties with the degree heuristic\\
    for value in DOMAIN(X) ordered by LCV:\\
    \quad try X=value
  \end{block}

  \begin{itemize}
    \item None of these heuristics removes a legal solution from the search space.
    \item They only alter the order in which equivalent search alternatives are explored.
    \item Their effectiveness comes from reaching contradictions or promising assignments sooner.
  \end{itemize}

  \begin{keyidea}
    Ordering helps us search more intelligently; inference goes further by deleting values that no longer need to be searched at all.
  \end{keyidea}
\end{frame}

\sectionframe{Part 3: Forward Checking --- Infer Immediate Consequences}

\begin{frame}[shrink=20]{Forward Checking}
  \begin{cpdefinition}
    \textbf{Forward checking} propagates each new assignment to its unassigned neighbors by deleting neighbor values inconsistent with that assignment.
  \end{cpdefinition}

  After assigning $X=v$:
  \begin{enumerate}
    \item inspect each unassigned neighbor $Y$ of $X$;
    \item remove every $y\in D_Y$ for which the constraint $C_{XY}(v,y)$ is false;
    \item fail immediately if any neighbor domain becomes empty.
  \end{enumerate}

  \begin{keyidea}
    A constraint violation need not wait until the neighbor is assigned; forward checking can predict the violation from the domain itself.
  \end{keyidea}
\end{frame}

\begin{frame}[shrink=20]{Forward Checking on Australia}
  \begin{columns}[T,onlytextwidth]
    \begin{column}{0.46\textwidth}
      Initially:
      \[
        D_{WA}=D_{NT}=D_{SA}=\{R,G,B\}.
      \]

      Assign
      \[
        WA=R.
      \]

      Because
      \[
        WA\neq NT,\qquad WA\neq SA,
      \]
      forward checking updates:
      \[
        D_{NT}\leftarrow\{G,B\},
      \]
      \[
        D_{SA}\leftarrow\{G,B\}.
      \]
    \end{column}

    \begin{column}{0.48\textwidth}
      \lectureimage[2.75in]{slide18.png}
    \end{column}
  \end{columns}
\end{frame}

\begin{frame}[shrink=20]{Failure Can Be Detected Before the Next Assignment}
  \begin{columns}[T,onlytextwidth]
    \begin{column}{0.52\textwidth}
      Suppose a neighbor currently has
      \[
        D_Y=\{R\}
      \]
      and a newly assigned adjacent variable becomes
      \[
        X=R
      \]
      under constraint $X\neq Y$.

      Forward checking removes $R$ from $D_Y$:
      \[
        D_Y\leftarrow\emptyset.
      \]

      The current branch is impossible even though $Y$ has not yet been assigned.
    \end{column}

    \begin{column}{0.42\textwidth}
      \lectureimage[2.50in]{slide19.png}
    \end{column}
  \end{columns}

  \begin{keyidea}
    An empty domain is a proof that the current partial assignment cannot be extended to a solution.
  \end{keyidea}
\end{frame}

\begin{frame}[shrink=20]{Forward Checking Inside Backtracking}
  \begin{block}{Algorithm sketch}
  \ttfamily\small
  for value in ORDERED-DOMAIN-VALUES(X):\\
  \quad if value is consistent with assignment:\\
  \qquad assign X=value\\
  \qquad pruned $\leftarrow$ FORWARD-CHECK(X, value)\\[0.03in]
  \qquad if no domain is empty:\\
  \qquad\quad result $\leftarrow$ BACKTRACK(assignment)\\
  \qquad\quad if result $\neq$ failure: return result\\[0.03in]
  \qquad RESTORE(pruned)\\
  \qquad remove X=value
  \end{block}

  \begin{keyidea}
    Inference changes the mutable domains, so every recursive branch must be able to undo exactly the values that branch pruned.
  \end{keyidea}
\end{frame}

\begin{frame}[shrink=20]{Implementation Detail: Domain Restoration}
  \begin{columns}[T,onlytextwidth]
    \begin{column}{0.49\textwidth}
      \begin{block}{Unsafe approach}
        Permanently delete inconsistent values from shared domain sets.

        After backtracking, later branches inherit pruning that was justified only by the abandoned assignment.
      \end{block}
    \end{column}

    \begin{column}{0.47\textwidth}
      \begin{block}{Safe approaches}
        \begin{itemize}
          \item Record each removed $(variable,value)$ pair and restore it.
          \item Copy domains for each recursive call.
          \item Use a reversible trail / stack of domain changes.
        \end{itemize}
      \end{block}
    \end{column}
  \end{columns}

  \begin{keyidea}
    Correct CSP implementations distinguish global problem constraints from temporary domain reductions caused by the current search branch.
  \end{keyidea}
\end{frame}

\sectionframe{Part 4: Arc Consistency and Constraint Propagation}

\begin{frame}[shrink=20]{Forward Checking Stops After One Step}
  \begin{columns}[T,onlytextwidth]
    \begin{column}{0.54\textwidth}
      Forward checking only propagates directly from the variable that was just assigned.

      Suppose:
      \[
        X\rightarrow Y\rightarrow Z.
      \]

      Assigning $X$ may shrink $D_Y$.

      That smaller $D_Y$ may make some values in $D_Z$ impossible, even though $Z$ is not adjacent to the newly assigned variable $X$.
    \end{column}

    \begin{column}{0.40\textwidth}
      \lectureimage[2.55in]{slide23.png}
    \end{column}
  \end{columns}

  \begin{keyidea}
    Constraint propagation repeatedly applies local reasoning until no additional domain values can be removed.
  \end{keyidea}
\end{frame}

\begin{frame}[shrink=20]{Arc Consistency --- Formal Definition}
  \begin{cpdefinition}
    A directed arc $X_i\rightarrow X_j$ is \textbf{arc-consistent} if every value $x\in D_i$ has at least one value $y\in D_j$ such that the binary constraint $C_{ij}(x,y)$ is satisfied.
  \end{cpdefinition}

  \[
    \forall x\in D_i,\ \exists y\in D_j:\ C_{ij}(x,y)=\mathrm{true}.
  \]

  \begin{itemize}
    \item The condition is directional: $X_i\rightarrow X_j$ and $X_j\rightarrow X_i$ are separate arcs.
    \item A value with no support in the neighboring domain can be deleted.
    \item Domain deletions may invalidate support on other arcs.
  \end{itemize}
\end{frame}

\begin{frame}[shrink=20]{Support Example}
  \begin{columns}[T,onlytextwidth]
    \begin{column}{0.50\textwidth}
      Let
      \[
        D_X=\{1,2,3\},\qquad D_Y=\{2,3\}
      \]
      with constraint
      \[
        X<Y.
      \]

      For $X\rightarrow Y$:
      \begin{itemize}
        \item $1$ is supported by $2$ or $3$;
        \item $2$ is supported by $3$;
        \item $3$ has no support.
      \end{itemize}

      Therefore remove $3$ from $D_X$.
    \end{column}

    \begin{column}{0.44\textwidth}
      \lectureimage[2.60in]{slide25.png}
    \end{column}
  \end{columns}
\end{frame}

\begin{frame}[shrink=20]{The REVISE Operation}
  \begin{block}{REVISE($X_i,X_j$)}
  \ttfamily\small
  revised $\leftarrow$ false\\
  for each $x$ in $D_i$:\\
  \quad if no $y$ in $D_j$ satisfies $C_{ij}(x,y)$:\\
  \qquad remove $x$ from $D_i$\\
  \qquad revised $\leftarrow$ true\\
  return revised
  \end{block}

  \begin{itemize}
    \item \texttt{REVISE} changes only $D_i$.
    \item It answers whether $D_i$ changed, because a change may affect other neighboring arcs.
    \item If $D_i$ becomes empty, the current CSP state is inconsistent.
  \end{itemize}
\end{frame}

\begin{frame}[shrink=20]{AC-3: Propagate Until Stable}
  \begin{cpdefinition}
    \textbf{AC-3} repeatedly revises directed arcs until every arc is consistent or some domain becomes empty.
  \end{cpdefinition}

  \begin{block}{High-level algorithm}
  \ttfamily\small
  queue $\leftarrow$ all directed arcs in the CSP\\[0.03in]
  while queue is not empty:\\
  \quad $(X_i,X_j) \leftarrow$ queue.pop()\\
  \quad if REVISE($X_i,X_j$):\\
  \qquad if $D_i$ is empty: return failure\\
  \qquad for each $X_k$ in NEIGHBORS($X_i$) except $X_j$:\\
  \qquad\quad queue.add(($X_k,X_i$))\\
  return success
  \end{block}
\end{frame}

\begin{frame}[shrink=20]{Why Re-Add Neighboring Arcs?}
  \begin{columns}[T,onlytextwidth]
    \begin{column}{0.54\textwidth}
      Suppose revising $Y\rightarrow Z$ removes a value from $D_Y$.

      A value in $D_X$ that previously depended on that deleted $Y$ value may now have no support.

      Therefore the arc
      \[
        X\rightarrow Y
      \]
      must be reconsidered.

      \begin{block}{Fixed-point view}
        AC-3 continues until every remaining value has local support and no revision can delete anything else.
      \end{block}
    \end{column}

    \begin{column}{0.40\textwidth}
      \lectureimage[2.55in]{slide28.png}
    \end{column}
  \end{columns}
\end{frame}

\begin{frame}[shrink=20]{AC-3 Trace: Initial Domains}
  \begin{columns}[T,onlytextwidth]
    \begin{column}{0.51\textwidth}
      Consider three variables:
      \[
        D_X=\{1,2,3\},\quad D_Y=\{1,2,3\},\quad D_Z=\{2\}
      \]
      with constraints
      \[
        X<Y,\qquad Y<Z.
      \]

      Start with arcs:
      \[
        X\rightarrow Y,
        Y\rightarrow X,
        Y\rightarrow Z,
        Z\rightarrow Y.
      \]

      Revising $Y\rightarrow Z$ leaves only
      \[
        D_Y=\{1\}.
      \]
    \end{column}

    \begin{column}{0.43\textwidth}
      \lectureplaceholder[2.75in]{Three-node chain X<Y<Z with domains shown; highlight Z={2} and the revision Y->Z pruning Y to {1}.}
    \end{column}
  \end{columns}
\end{frame}

\begin{frame}[shrink=20]{AC-3 Trace: Propagation Reveals Failure}
  With
  \[
    D_Y=\{1\},
  \]
  reconsider the neighboring arc
  \[
    X\rightarrow Y.
  \]

  The constraint is $X<Y$, but no value in
  \[
    D_X=\{1,2,3\}
  \]
  is less than $1$.

  Therefore
  \[
    D_X\leftarrow\emptyset.
  \]

  \begin{keyidea}
    AC-3 derives a contradiction without assigning $X$ or $Y$ in search; the existing domains and constraints are already sufficient to prove failure.
  \end{keyidea}
\end{frame}

\begin{frame}[shrink=20]{AC-3 Cost and Benefit}
  \begin{columns}[T,onlytextwidth]
    \begin{column}{0.48\textwidth}
      \begin{block}{Cost}
        For a binary CSP with at most $d$ values per domain and $e$ directed arcs, a straightforward AC-3 implementation has worst-case time on the order of
        \[
          O(ed^3).
        \]
      \end{block}
    \end{column}

    \begin{column}{0.48\textwidth}
      \begin{block}{Benefit}
        \begin{itemize}
          \item deletes values before branching;
          \item can expose contradictions early;
          \item may dramatically reduce the later backtracking tree.
        \end{itemize}
      \end{block}
    \end{column}
  \end{columns}

  \begin{keyidea}
    Constraint propagation deliberately spends computation now to avoid potentially exponential search later.
  \end{keyidea}
\end{frame}

\begin{frame}[shrink=20]{Forward Checking vs. Arc Consistency}
  \begin{center}
  \renewcommand{\arraystretch}{1.35}
  \begin{tabular}{l|l|l}
    & \textbf{Forward Checking} & \textbf{Arc Consistency} \\\hline
    Trigger & new assignment & domain changes / queued arcs \\
    Propagation & assigned variable to neighbors & repeatedly across neighboring arcs \\
    Cost & lower & higher \\
    Pruning strength & weaker & stronger \\
    Indirect contradictions & may miss & may expose
  \end{tabular}
  \end{center}

  \begin{keyidea}
    Forward checking looks one step ahead; arc consistency repeatedly propagates the consequences of shrinking domains.
  \end{keyidea}
\end{frame}

\begin{frame}[shrink=20]{Maintaining Arc Consistency (MAC)}
  \begin{cpdefinition}
    \textbf{Maintaining Arc Consistency (MAC)} runs arc-consistency propagation after each tentative assignment during backtracking search.
  \end{cpdefinition}

  \begin{enumerate}
    \item Choose a variable and tentative value.
    \item Restrict that variable's domain to the chosen value.
    \item Run AC-style propagation on affected arcs.
    \item Recurse only if every domain remains nonempty.
    \item Restore pruned values when backtracking.
  \end{enumerate}

  \begin{keyidea}
    MAC tightly couples search and inference: every decision is immediately followed by the strongest local consequences the solver knows how to derive.
  \end{keyidea}
\end{frame}

\begin{frame}[shrink=20]{Backtracking With Heuristics and MAC}
  \begin{block}{Integrated solver}
  \ttfamily\small
  BACKTRACK(assignment, domains):\\
  \quad if assignment is complete: return assignment\\
  \quad X $\leftarrow$ SELECT-VARIABLE(MRV, degree)\\[0.03in]
  \quad for value in ORDER-VALUES(X, LCV):\\
  \qquad if value is consistent:\\
  \qquad\quad assign X=value\\
  \qquad\quad restrict $D_X$ to \{value\}\\
  \qquad\quad if MAC(domains) succeeds:\\
  \qquad\qquad result $\leftarrow$ BACKTRACK(...)\\
  \qquad\qquad if result $\neq$ failure: return result\\
  \qquad\quad restore domains and assignment\\
  \quad return failure
  \end{block}

  \begin{keyidea}
    A practical CSP solver is a composition of search order, reversible domain state, and inference.
  \end{keyidea}
\end{frame}

\sectionframe{Part 5: Recovering More Intelligently From Failure}

\begin{frame}[shrink=20]{Chronological Backtracking Can Undo the Wrong Decision}
  \begin{columns}[T,onlytextwidth]
    \begin{column}{0.54\textwidth}
      Standard backtracking returns to the most recently assigned variable.

      But suppose a contradiction is caused by assignments to
      \[
        X_2\quad\text{and}\quad X_7,
      \]
      while $X_6$ has no relationship to the failed constraint.

      Chronological backtracking may first reconsider $X_6$ even though changing it cannot repair the conflict.
    \end{column}

    \begin{column}{0.40\textwidth}
      \lectureimage[2.55in]{slide36.png}
    \end{column}
  \end{columns}
\end{frame}

\begin{frame}[shrink=20]{Backjumping}
  \begin{cpdefinition}
    \textbf{Backjumping} uses information about which earlier assignments contributed to a failure and jumps back to a relevant conflicting variable rather than always returning one level.
  \end{cpdefinition}

  \begin{itemize}
    \item The solver records or reconstructs a conflict set associated with failure.
    \item Variables unrelated to the conflict can be skipped.
    \item More advanced variants combine conflict information across failures.
  \end{itemize}

  \begin{keyidea}
    Ordinary backtracking asks, ``what did I choose most recently?'' Backjumping asks, ``which choice actually caused this contradiction?''
  \end{keyidea}
\end{frame}

\sectionframe{Part 6: Min-Conflicts --- Local Search for Constraint Satisfaction}

\begin{frame}[shrink=20]{A Different CSP Strategy}
  \begin{columns}[T,onlytextwidth]
    \begin{column}{0.47\textwidth}
      \begin{block}{Backtracking family}
        \begin{itemize}
          \item starts with a partial assignment;
          \item maintains consistency while adding values;
          \item can systematically prove failure.
        \end{itemize}
      \end{block}
    \end{column}

    \begin{column}{0.47\textwidth}
      \begin{block}{Local-search family}
        \begin{itemize}
          \item starts with a complete assignment;
          \item allows temporary constraint violations;
          \item repeatedly repairs conflicts.
        \end{itemize}
      \end{block}
    \end{column}
  \end{columns}

  \begin{keyidea}
    CSPs can be solved either by constructing a legal assignment or by repairing an illegal complete assignment.
  \end{keyidea}
\end{frame}

\begin{frame}[shrink=20]{Min-Conflicts}
  \begin{cpdefinition}
    \textbf{Min-conflicts} is a local-search method that repeatedly selects a conflicted variable and assigns the value that minimizes the number of violated constraints.
  \end{cpdefinition}

  \begin{block}{Algorithm}
  \ttfamily\small
  current $\leftarrow$ a complete assignment\\
  repeat up to max\_steps:\\
  \quad if current satisfies every constraint: return current\\
  \quad X $\leftarrow$ randomly choose a conflicted variable\\
  \quad value $\leftarrow$ value minimizing CONFLICTS(X, value, current)\\
  \quad set X=value\\
  return failure
  \end{block}
\end{frame}

\begin{frame}[shrink=20]{Min-Conflicts on N-Queens}
  \begin{columns}[T,onlytextwidth]
    \begin{column}{0.53\textwidth}
      Use one variable per column:
      \[
        Q_1,Q_2,\ldots,Q_n.
      \]

      Domain:
      \[
        D_i=\{1,2,\ldots,n\}
      \]
      gives the row position of queen $i$.

      At each step:
      \begin{enumerate}
        \item choose a queen currently attacking another queen;
        \item move it to the row that creates the fewest attacks;
        \item repeat until no conflicts remain.
      \end{enumerate}
    \end{column}

    \begin{column}{0.41\textwidth}
      \lectureimage[2.70in]{slide41.png}
    \end{column}
  \end{columns}
\end{frame}

\begin{frame}[shrink=20]{When Min-Conflicts Is Attractive}
  \begin{columns}[T,onlytextwidth]
    \begin{column}{0.47\textwidth}
      \begin{block}{Strengths}
        \begin{itemize}
          \item very little memory;
          \item simple local updates;
          \item often effective when solutions are plentiful;
          \item natural for large assignment problems with easy repair moves.
        \end{itemize}
      \end{block}
    \end{column}

    \begin{column}{0.47\textwidth}
      \begin{block}{Limitations}
        \begin{itemize}
          \item incomplete under a finite step limit;
          \item may cycle or become trapped;
          \item does not naturally prove unsatisfiability;
          \item performance depends on landscape structure and restart strategy.
        \end{itemize}
      \end{block}
    \end{column}
  \end{columns}

  \begin{keyidea}
    Min-conflicts trades systematic guarantees for the ability to move through very large complete-state spaces cheaply.
  \end{keyidea}
\end{frame}

\begin{frame}[shrink=20]{Sudoku: Systematic Search or Min-Conflicts?}
  \begin{columns}[T,onlytextwidth]
    \begin{column}{0.48\textwidth}
      \begin{block}{Backtracking + propagation}
        \begin{itemize}
          \item exploits forced values aggressively;
          \item naturally handles tightly constrained puzzles;
          \item can establish failure when no solution exists.
        \end{itemize}
      \end{block}
    \end{column}

    \begin{column}{0.48\textwidth}
      \begin{block}{Min-conflicts}
        \begin{itemize}
          \item begins from a complete grid assignment;
          \item repairs row, column, and box conflicts;
          \item may struggle when legal configurations occupy narrow regions of the state space.
        \end{itemize}
      \end{block}
    \end{column}
  \end{columns}

  \begin{keyidea}
    The best CSP algorithm depends on whether the problem rewards logical propagation, local repair, or some combination of both.
  \end{keyidea}
\end{frame}

\sectionframe{Part 7: Applications Beyond Static Puzzles}

\begin{frame}[shrink=20]{Scheduling as a CSP}
  \begin{columns}[T,onlytextwidth]
    \begin{column}{0.55\textwidth}
      A scheduling problem can assign each activity a time, room, or resource.

      \begin{block}{Variables}
        Courses, jobs, meetings, or tasks.
      \end{block}

      \begin{block}{Domains}
        Feasible times, rooms, machines, or workers.
      \end{block}

      \begin{block}{Constraints}
        No overlap, precedence, capacity, availability, and resource limits.
      \end{block}
    \end{column}

    \begin{column}{0.39\textwidth}
      \lectureplaceholder[2.65in]{Timetable or job-shop diagram with overlapping assignments marked invalid and precedence arrows between selected tasks.}
    \end{column}
  \end{columns}

  \begin{keyidea}
    The same heuristics and propagation mechanisms used for maps and Sudoku also apply to real allocation and scheduling systems.
  \end{keyidea}
\end{frame}

\begin{frame}[shrink=20]{What If Constraints Are Revealed Over Time?}
  \begin{columns}[T,onlytextwidth]
    \begin{column}{0.54\textwidth}
      Classical CSPs assume the variables, domains, and constraints are known before solving.

      Some interactive problems instead reveal information after each action:
      \begin{itemize}
        \item diagnosis through tests;
        \item preference elicitation;
        \item planning with feedback;
        \item word-guessing games.
      \end{itemize}

      The solver repeatedly updates domains and constraints as new observations arrive.
    \end{column}

    \begin{column}{0.40\textwidth}
      \lectureplaceholder[2.55in]{Interactive loop: choose action/guess -> receive feedback -> update domains/constraints -> choose again.}
    \end{column}
  \end{columns}
\end{frame}

\begin{frame}[shrink=20]{Wordle as Constraint Accumulation}
  \begin{columns}[T,onlytextwidth]
    \begin{column}{0.52\textwidth}
      Treat the five letter positions as variables with letter domains.

      Feedback introduces constraints:
      \begin{itemize}
        \item \textbf{green}: the letter is fixed at that position;
        \item \textbf{yellow}: the letter occurs elsewhere but not at that position;
        \item \textbf{gray}: the guessed occurrence is excluded, subject to repeated-letter counts.
      \end{itemize}

      Each guess filters the remaining candidate assignments.
    \end{column}

    \begin{column}{0.42\textwidth}
      \lectureplaceholder[2.70in]{Wordle-style five-letter grid with green, yellow, and gray feedback, alongside shrinking candidate letter domains.}
    \end{column}
  \end{columns}

  \begin{keyidea}
    Constraint reasoning remains useful even when the complete constraint set is discovered incrementally rather than given all at once.
  \end{keyidea}
\end{frame}

\sectionframe{Part 8: Putting the Solver Together}

\begin{frame}[shrink=20]{Architecture of a Practical Backtracking Solver}
  \begin{center}
  \lectureplaceholder[3.05in]{Pipeline diagram: BACKTRACK -> SELECT VARIABLE (MRV + degree) -> ORDER VALUES (LCV) -> ASSIGN -> INFERENCE (forward checking or MAC/AC-3) -> recurse / restore on failure.}
  \end{center}
\end{frame}

\begin{frame}[shrink=20]{What Each Technique Contributes}
  \begin{center}
  \renewcommand{\arraystretch}{1.30}
  \begin{tabular}{l|l}
    \textbf{Technique} & \textbf{Main question answered} \\\hline
    MRV & Which variable is closest to failure? \\
    Degree heuristic & Which variable constrains the most remaining variables? \\
    LCV & Which value preserves the most flexibility? \\
    Forward checking & What immediate neighbor values became impossible? \\
    AC-3 & What additional impossibilities propagate through the graph? \\
    MAC & How do we maintain propagation throughout search? \\
    Backjumping & Which earlier decision actually caused this failure? \\
    Min-conflicts & Can a complete assignment be repaired instead?
  \end{tabular}
  \end{center}
\end{frame}

\begin{frame}[shrink=20]{Choosing a CSP Strategy}
  \begin{columns}[T,onlytextwidth]
    \begin{column}{0.47\textwidth}
      \begin{block}{Prefer systematic search when}
        \begin{itemize}
          \item completeness matters;
          \item proving unsatisfiability matters;
          \item strong propagation is available;
          \item constraints make many values forced or impossible.
        \end{itemize}
      \end{block}
    \end{column}

    \begin{column}{0.47\textwidth}
      \begin{block}{Prefer local repair when}
        \begin{itemize}
          \item the state space is enormous;
          \item complete assignments are easy to construct;
          \item conflicts are cheap to evaluate and repair;
          \item finding a good solution matters more than proving failure.
        \end{itemize}
      \end{block}
    \end{column}
  \end{columns}

  \begin{keyidea}
    Constraint satisfaction is a modeling framework; the solving strategy should match the guarantees and structure the application actually requires.
  \end{keyidea}
\end{frame}

\begin{frame}[shrink=20]{Constraint Satisfaction I and II}
  \begin{center}
  \renewcommand{\arraystretch}{1.32}
  \begin{tabular}{l|l}
    \textbf{Lecture 8: Foundation} & \textbf{Lecture 9: Efficient Solving} \\\hline
    Variables, domains, constraints & MRV, degree, LCV \\
    Partial and complete assignments & Domain-aware ordering \\
    Constraint graphs & Forward checking \\
    Consistency intuition & Arc consistency and AC-3 \\
    Basic backtracking & MAC and reversible inference \\
    Map coloring and Sudoku modeling & Backjumping and min-conflicts
  \end{tabular}
  \end{center}

  \begin{keyidea}
    The CSP representation becomes powerful when the solver uses constraint structure not only to test assignments, but to decide what to search and what can be inferred without search.
  \end{keyidea}
\end{frame}

\begin{frame}[shrink=20]{Takeaways}
  \begin{itemize}
    \item MRV and the degree heuristic choose informative variables; LCV preserves flexibility when choosing values.
    \item Forward checking removes values made immediately inconsistent by a new assignment.
    \item Arc consistency requires every value to have support in each neighboring domain.
    \item AC-3 repeatedly applies \texttt{REVISE} until a fixed point or an empty domain is reached.
    \item MAC embeds arc-consistency propagation inside backtracking and requires reversible domain updates.
    \item Backjumping uses conflict information to avoid undoing irrelevant decisions.
    \item Min-conflicts solves CSPs from the opposite direction by repairing complete but inconsistent assignments.
    \item Effective CSP solving combines representation, search order, inference, and problem-specific structure.
  \end{itemize}

  \begin{keyidea}
    The central algorithmic lesson of CSPs is that reasoning can remove large regions of a combinatorial search space before those regions are ever explored.
  \end{keyidea}
\end{frame}

\end{document}
